% Generated draft: requires mathematical and source review.
\documentclass[11pt]{article}
\usepackage[margin=25mm]{geometry}
\usepackage{fontspec}
\setmainfont{Latin Modern Roman}
\usepackage{amsmath,amssymb,mathtools,unicode-math}
\setmathfont{Latin Modern Math}
\usepackage{enumitem,needspace,xurl,hyperref}
\setlength{\parindent}{0pt}
\setlength{\parskip}{5pt}
\setlength{\emergencystretch}{3em}
\newcommand{\N}{\mathbb{N}}
\newcommand{\Z}{\mathbb{Z}}
\newcommand{\Q}{\mathbb{Q}}
\newcommand{\R}{\mathbb{R}}
\newcommand{\C}{\mathbb{C}}
\newcommand{\F}{\mathbb{F}}
\newcommand{\A}{\mathbb{A}}
\newcommand{\PP}{\mathbb{P}}
\newcommand{\E}{\mathbb{E}}
\newcommand{\eps}{\varepsilon}
\newcommand{\dd}{\,\mathrm{d}}
\begin{document}
\section*{UnsolvedMath ID: 30003360}
\textbf{Source problem code:} catalogue entry
\subsection{Definitions and mathematical statement}
Let \(M\) denote the underlying smooth manifold. The source does not specify its dimension, compactness, boundary conditions, or topology. Write
\[
\operatorname{Riem}(M)=\{g\in\Gamma^{\infty}(\operatorname{Sym}^{2}T^{*}M):g_x(v,v)>0\text{ for every }x\in M\text{ and }0\ne v\in T_xM\}
\]
for its space of smooth Riemannian metrics, and let \(\operatorname{Diff}(M)\) be its group of smooth diffeomorphisms. Pullback defines a right action \(g\cdot\varphi=\varphi^{*}g\). The stabilizer of a metric is
\[
\operatorname{Stab}(g)=\{\varphi\in\operatorname{Diff}(M):\varphi^{*}g=g\},
\]
and the set of orbits is
\[
\mathcal M(M)=\operatorname{Riem}(M)/\operatorname{Diff}(M).
\]
The action is nonfree precisely when there exist \(g\in\operatorname{Riem}(M)\) and \(\varphi\ne\operatorname{id}_M\) with \(\varphi^{*}g=g\). These definitions describe the unrestricted smooth metric space mentioned in the statement; the source supplies no topology or analytic completion for the metric space or the acting group.

The supplied question is: in a setting where this action is nonfree and the resulting moduli space has what the source calls ``orbifold singularities,'' what mathematical framework or construction should be used to analyze that moduli space?

The presence of orbifold singularities must be retained as a separate source qualification: nonfreeness alone does not establish an orbifold structure. The record specifies neither local quotient models nor hypotheses on stabilizers that would justify that terminology. It also supplies no precise property that the requested framework must establish. Thus this is a methodological question with unresolved hypotheses and an unspecified mathematical target, rather than a fully quantified conjecture. No universal quantifier over manifolds, scalar curvature restriction, or replacement of smooth metrics by singular metrics is imposed by the supplied statement.
\subsection{Short English statement}
The question asks how to study a moduli space of Riemannian metrics when some metrics have nontrivial diffeomorphism symmetries and the quotient is described as having orbifold singularities. The source leaves the manifold, the analytic setting, the meaning of those singularities, and the desired outcome unspecified. Its separate literature assessment concerns different extension questions and does not resolve these omissions.
\subsection{Sources}
\begin{itemize}
\item Supplied local UnsolvedMath registry record, permanent ID \(30003360\), source code \texttt{catalogue entry}, title ``Orbifold Singularities in Riemannian Metric Moduli Spaces.'' Project registry locations: \texttt{data/unsolvedmath/problems.json} and \texttt{catalog/problems.json}; this draft uses the supplied complete record without independently inspecting those files. The record labels the problem open and its statement exact, with verification described as internal readability and byte-for-byte retention. These are supplied curation labels, not independent mathematical or source verification.
\item Supplied source citation: \emph{Mini-Workshop: Spaces and Moduli Spaces of Riemannian Metrics (2018)}, Oberwolfach Reports, pp.\ 133--166, DOI \texttt{10.4171/owr/2017/3}. \url{https://doi.org/10.4171/owr/2017/3}. The background describes this as a workshop report documenting mathematical discussions and open problems. The registry gives proposed year \(2017\), whereas the supplied citation includes \(2018\); this discrepancy is unresolved. No individual proposer is supplied.
\item Supplied literature source, ``Oberwolfach discussion of open problems'': \url{https://doi.org/10.4171/owr/2017/3}.
\item Supplied literature source, ``Positive scalar curvature on manifolds with singularities'': \url{https://doi.org/10.3842/SIGMA.2021.034}.
\item Literature assessment dated exactly \texttt{2026-08-21}: supplied dated curation, not independently verified. It labels the problem open and states that the record combines questions about extending eta-invariant results to nonnegative scalar curvature and replacing smooth metrics by singular metrics, with neither general extension established. The background identifies this as dated literature triage derived from supplied Oberwolfach curation and requests independent verification of time-sensitive status and cited claims. Those extension questions do not match the displayed question about nonfree diffeomorphism actions; their relevance remains unresolved.
\end{itemize}
\end{document}
